\subsection{Homological evaluation and intrinsic visibility}

The strict quotient $A_{\mathrm{vis}}=A/H_\alpha$ depends on a chosen cocycle representative. To isolate the part of the obstruction that belongs to the class itself, we now make the homological evaluation map primary. The intrinsic quotient will be defined from this map and only afterwards recovered through class-admissible character channels.

\begin{definition}[Finite nerve presentation of the obstruction class]\label{def:finite-nerve-presentation}
Fix an admitted reference $r:\mathsf{Ref}_s$ at a state $p$, and assume conditions \textup{(i)}--\textup{(iv)} of \Cref{def:finite-obstruction-presentation}. A \emph{finite nerve presentation} of the branch obstruction class consists of:
\begin{enumerate}[label=(\roman*),nosep]
\item a finite covering family
\[
\mathcal U=\{a_i\to a\}_{i\in I},
\qquad
a=[r]_p,
\]
whose Cech nerve $N(\mathcal U)$ is finite;
\item a trivialization of the band on $\mathcal U$, so that the band may be viewed as a constant coefficient finite abelian group $A$ on $N(\mathcal U)$;
\item the induced obstruction class
\[
\omega\in H^2(N(\mathcal U),A),
\]
obtained by identifying Cech cochains on $\mathcal U$ with simplicial cochains on the finite nerve and writing
\[
\omega:=\omega_v
\]
for the chosen visible branch.
\end{enumerate}
Write
\[
\eta_A:H^2(N(\mathcal U),A)\longrightarrow \mathrm{Hom}(H_2(N(\mathcal U),\mathbb Z),A)
\]
for the map in the universal coefficient exact sequence
\[
0\to \operatorname{Ext}^1(H_1(N(\mathcal U),\mathbb Z),A)\to H^2(N(\mathcal U),A)\xrightarrow{\eta_A}\mathrm{Hom}(H_2(N(\mathcal U),\mathbb Z),A)\to 0
\]
\cite[Sec.~3.1]{Hatcher2002}, and define the \emph{homological evaluation map}
\[
\mathrm{ev}_{\omega}:=\eta_A(\omega).
\]
\end{definition}

\begin{definition}[Intrinsic visible quotient]\label{def:intrinsic-visible-quotient}
In the setting of \Cref{def:finite-nerve-presentation}, define the \emph{intrinsic hidden subgroup}
\[
K_{\omega}:=\mathrm{Im}(\mathrm{ev}_{\omega})
\]
and the associated \emph{intrinsic visible quotient}
\[
A_{\mathrm{vis}}^{\omega}:=\operatorname{coker}(\mathrm{ev}_{\omega})=A/K_{\omega}.
\]
\end{definition}

\begin{definition}[Class-admissible characters]\label{def:class-admissible-character}
In the setting of \Cref{def:finite-nerve-presentation}, let
\[
X_{\omega}:=\{\chi\in\widehat A\mid \chi_*(\omega)=0\},
\qquad
\widehat A:=\mathrm{Hom}(A,\mathbb T).
\]
Elements of $X_{\omega}$ are called \emph{class-admissible} for $\omega$.
\end{definition}

\begin{theorem}[Character-theoretic detection of the intrinsic visible quotient]\label{thm:intrinsic-visible-quotient}
In the setting of \Cref{def:class-admissible-character},
\[
X_{\omega}
=
\{\chi\in\widehat A\mid \chi\circ \mathrm{ev}_{\omega}=0\}
=
\mathrm{Ann}(K_{\omega}).
\]
Consequently,
\[
K_{\omega}
=
\bigcap_{\chi\in X_{\omega}}\ker\chi,
\qquad
A_{\mathrm{vis}}^{\omega}\cong A\Big/\!\bigcap_{\chi\in X_{\omega}}\ker\chi.
\]
\end{theorem}

\begin{proof}
For every character $\chi:A\to\mathbb T$, naturality of the universal coefficient sequence gives a commutative square
\[
\begin{CD}
H^2(N(\mathcal U),A) @>\eta_A>> \mathrm{Hom}(H_2(N(\mathcal U),\mathbb Z),A) \\
@V\chi_*VV @VV(\chi\circ -)V \\
H^2(N(\mathcal U),\mathbb T) @>\eta_{\mathbb T}>> \mathrm{Hom}(H_2(N(\mathcal U),\mathbb Z),\mathbb T).
\end{CD}
\]
Since $\mathbb T$ is divisible and hence injective in $\mathbf{Ab}$, one has
\[
\operatorname{Ext}^1(H_1(N(\mathcal U),\mathbb Z),\mathbb T)=0
\]
\cite[Sec.~2.3]{Weibel1994}. Thus $\eta_{\mathbb T}$ is an isomorphism by \Cref{def:finite-nerve-presentation}. Therefore
\[
\chi_*(\omega)=0
\Longleftrightarrow
\eta_{\mathbb T}(\chi_*(\omega))=0
\Longleftrightarrow
\chi\circ \eta_A(\omega)=0
\Longleftrightarrow
\chi\circ \mathrm{ev}_{\omega}=0.
\]
This proves the first displayed identity, and the second is simply the definition of the annihilator of a subgroup of $A$.

Let $H:=K_{\omega}=\mathrm{Im}(\mathrm{ev}_{\omega})\le A$. By the first part of the proof, $X_{\omega}=\mathrm{Ann}(H)$. Finite abelian duality implies
\[
\bigcap_{\chi\in \mathrm{Ann}(H)}\ker\chi=H
\]
\cite[Chaps.~5, 12]{DummitFoote2004}. Hence
\[
K_{\omega}
=
\bigcap_{\chi\in X_{\omega}}\ker\chi.
\]
The final isomorphism follows immediately:
\[
A_{\mathrm{vis}}^{\omega}
=
A/K_{\omega}
\cong
A\Big/\!\bigcap_{\chi\in X_{\omega}}\ker\chi.
\qedhere
\]
\end{proof}

\begin{corollary}[Character detection of the intrinsic visible quotient]\label{cor:intrinsic-character-detection}
The quotient map
\[
\pi_{\omega}:A\twoheadrightarrow A_{\mathrm{vis}}^{\omega}
\]
induces a natural isomorphism
\[
\widehat{A_{\mathrm{vis}}^{\omega}}\cong X_{\omega}.
\]
For elements $x,y\in A$, one has
\[
\pi_{\omega}(x)=\pi_{\omega}(y)
\quad\Longleftrightarrow\quad
\chi(x)=\chi(y)\ \text{for every class-admissible }\chi.
\]
\end{corollary}

\begin{proof}
By \Cref{thm:intrinsic-visible-quotient},
\[
X_{\omega}=\mathrm{Ann}(K_{\omega})=K_{\omega}^{\perp}.
\]
Characters of $A_{\mathrm{vis}}^{\omega}=A/K_{\omega}$ are exactly the characters of $A$ that kill $K_{\omega}$, so precomposition with $\pi_{\omega}$ identifies $\widehat{A_{\mathrm{vis}}^{\omega}}$ with $X_{\omega}$. The separation statement follows from finite abelian duality applied to the finite group $A_{\mathrm{vis}}^{\omega}$ \cite[Chap.~5]{DummitFoote2004}.
\end{proof}

\begin{theorem}[Initiality among pure-extension quotients]\label{thm:intrinsic-pure-ext-initiality}
Let
\[
q:A\twoheadrightarrow B
\]
be a surjective homomorphism onto a finite abelian group $B$. In the setting of \Cref{def:finite-nerve-presentation}, the following are equivalent:
\begin{enumerate}[label=(\roman*),nosep]
\item the pushed-forward class $q_*(\omega)\in H^2(N(\mathcal U),B)$ lies in the subgroup
\[
\operatorname{Ext}^1(H_1(N(\mathcal U),\mathbb Z),B);
\]
\item $\eta_B(q_*(\omega))=0$;
\item $q\circ \mathrm{ev}_{\omega}=0$;
\item $\mathrm{Im}(\mathrm{ev}_{\omega})\subseteq \ker q$;
\item $q$ factors uniquely through
\[
\pi_{\omega}:A\twoheadrightarrow A_{\mathrm{vis}}^{\omega}.
\]
\end{enumerate}
Hence $A_{\mathrm{vis}}^{\omega}$ is the initial coefficient quotient through which the $H_2$-visible part of the obstruction disappears, leaving only a pure $\operatorname{Ext}$-type remnant.
\end{theorem}

\begin{proof}
The equivalence of \textup{(i)} and \textup{(ii)} is exactly the universal coefficient sequence for $B$ \cite[Sec.~3.1]{Hatcher2002}. Naturality gives
\[
\eta_B(q_*(\omega))=q\circ \eta_A(\omega)=q\circ \mathrm{ev}_{\omega},
\]
so \textup{(ii)} and \textup{(iii)} are equivalent. Conditions \textup{(iii)} and \textup{(iv)} are equivalent by definition of the image of a homomorphism. By \Cref{def:intrinsic-visible-quotient}, $\mathrm{Im}(\mathrm{ev}_{\omega})=K_{\omega}$, so \textup{(iv)} is equivalent to
\[
K_{\omega}\subseteq \ker q.
\]
By the universal property of the quotient by $K_{\omega}$, this is equivalent to unique factorization of $q$ through $\pi_{\omega}$, proving \textup{(v)}.
\end{proof}

\begin{corollary}[Class-trivializing initiality when $H_1$ vanishes]\label{cor:h1-vanishing-class-trivializing}
Assume
\[
H_1(N(\mathcal U),\mathbb Z)=0.
\]
Then, for every surjective homomorphism $q:A\twoheadrightarrow B$, the following are equivalent:
\begin{enumerate}[label=(\roman*),nosep]
\item $q_*(\omega)=0$ in $H^2(N(\mathcal U),B)$;
\item $q$ factors uniquely through
\[
\pi_{\omega}:A\twoheadrightarrow A_{\mathrm{vis}}^{\omega}.
\]
\end{enumerate}
Hence, when $H_1(N(\mathcal U),\mathbb Z)=0$, the intrinsic visible quotient is initial among coefficient quotients that annihilate the obstruction class itself.
\end{corollary}

\begin{proof}
If $H_1(N(\mathcal U),\mathbb Z)=0$, then
\[
\operatorname{Ext}^1(H_1(N(\mathcal U),\mathbb Z),B)=0
\]
for every abelian group $B$. Thus the pure-extension condition in \Cref{thm:intrinsic-pure-ext-initiality}(i) is equivalent to $q_*(\omega)=0$, and the conclusion follows.
\end{proof}

\begin{remark}[Canonical visibility without a canonical splitting]\label{rem:no-canonical-splitting}
The quotient $A_{\mathrm{vis}}^{\omega}$ and the subgroup $K_{\omega}=\mathrm{Im}(\mathrm{ev}_{\omega})$ are canonical. By contrast, the universal coefficient sequence need not split canonically. Thus the theory isolates a canonical visible quotient and a canonical blind subgroup without requiring any noncanonical decomposition of $\omega$ into visible and invisible summands.
\end{remark}

\begin{corollary}[Comparison with strict visible quotients]\label{cor:strict-to-intrinsic-visible}
Let $\alpha$ be any finite obstruction presentation of the class $\omega$. Then the intrinsic visible quotient
\[
A_{\mathrm{vis}}^{\omega}
\]
admits a canonical surjective homomorphism onto the strict visible quotient
\[
A/H_\alpha.
\]
Equivalently, every strict visible quotient is a quotient of the intrinsic visible quotient.
\end{corollary}

\begin{proof}
Let $q:A\twoheadrightarrow A/H_\alpha$ be the canonical quotient. By \Cref{thm:initial-visible-quotient}, the quotient kills every cocycle value $\alpha_{ijk}$ pointwise, so the pushed-forward cocycle representative is zero. Hence $q_*(\omega)=0$, and in particular it lies in the $\operatorname{Ext}$-subgroup of $H^2(N(\mathcal U),A/H_\alpha)$. Apply \Cref{thm:intrinsic-pure-ext-initiality}.
\end{proof}

\begin{definition}[Chain-level evaluation of a presentation]\label{def:chain-level-evaluation}
Let $\alpha\in Z^2(N(\mathcal U),A)$ be a finite obstruction presentation of the class $\omega$. Write
\[
C_2(N(\mathcal U),\mathbb Z)
\]
for the simplicial group of integral $2$-chains on the finite nerve. The \emph{chain-level evaluation map} of $\alpha$ is the homomorphism
\[
\widetilde\alpha:C_2(N(\mathcal U),\mathbb Z)\to A
\]
defined by linear extension from the oriented $2$-simplices:
\[
\widetilde\alpha\!\left(\sum_{\sigma} n_{\sigma}\sigma\right)
:=
\sum_{\sigma} n_{\sigma}\alpha(\sigma).
\]
\end{definition}

\begin{proposition}[Strict visibility versus homological visibility]\label{prop:chain-vs-homological-visibility}
In the setting of \Cref{def:chain-level-evaluation}, the following hold.
\begin{enumerate}[label=(\roman*),nosep]
\item $\widetilde\alpha$ vanishes on $2$-boundaries and therefore induces a homomorphism
\[
\overline\alpha:H_2(N(\mathcal U),\mathbb Z)\to A.
\]
\item The induced homomorphism $\overline\alpha$ is exactly the homological evaluation map $\mathrm{ev}_{\omega}$.
\item The strict obstruction subgroup and the intrinsic hidden subgroup satisfy
\[
H_\alpha=\mathrm{Im}(\widetilde\alpha),
\qquad
K_{\omega}=\mathrm{Im}(\overline\alpha)=\widetilde\alpha(Z_2(N(\mathcal U),\mathbb Z)).
\]
\end{enumerate}
Thus the gap between strict visibility and intrinsic visibility is exactly the gap between the image of the cocycle on all $2$-chains and its image on $2$-cycles.
\end{proposition}

\begin{proof}
If $c\in C_3(N(\mathcal U),\mathbb Z)$, then
\[
\widetilde\alpha(\partial c)=(\delta\alpha)(c)=0
\]
because $\alpha$ is a cocycle. Hence $\widetilde\alpha$ vanishes on $B_2(N(\mathcal U),\mathbb Z)$ and induces a homomorphism
\[
\overline\alpha:H_2(N(\mathcal U),\mathbb Z)=Z_2/B_2\to A.
\]
This proves \textup{(i)}.

By definition, the cohomology class of $\alpha$ is $\omega$, and evaluation of a cocycle on cycles is precisely the map in the universal coefficient theorem associated with that class. Therefore $\overline\alpha=\eta_A(\omega)=\mathrm{ev}_{\omega}$, which is \textup{(ii)}.

For \textup{(iii)}, the image of $\widetilde\alpha$ is exactly the subgroup generated by the simplex values $\alpha_{ijk}$, hence it is $H_\alpha$ by \Cref{def:finite-obstruction-presentation}. The equality
\[
\mathrm{Im}(\overline\alpha)=\widetilde\alpha(Z_2(N(\mathcal U),\mathbb Z))
\]
is immediate from \textup{(i)}, and \Cref{def:intrinsic-visible-quotient} identifies this image with $K_{\omega}$.
\end{proof}

\begin{corollary}[Criterion for coincidence of strict and intrinsic visibility]\label{cor:strict-intrinsic-coincidence-criterion}
Let $\alpha$ be a finite obstruction presentation of $\omega$. The following are equivalent:
\begin{enumerate}[label=(\roman*),nosep]
\item $A/H_\alpha\cong A_{\mathrm{vis}}^{\omega}$;
\item $H_\alpha=K_{\omega}$;
\item $\widetilde\alpha(C_2(N(\mathcal U),\mathbb Z))=\widetilde\alpha(Z_2(N(\mathcal U),\mathbb Z))$.
\end{enumerate}
Equivalently, strict visibility agrees with intrinsic visibility exactly when every presentation-level visible value already arises from the evaluation of the cocycle on a genuine $2$-cycle.
\end{corollary}

\begin{proof}
The equivalence of \textup{(i)} and \textup{(ii)} is immediate from the quotient descriptions
\[
A/H_\alpha
\qquad\text{and}\qquad
A_{\mathrm{vis}}^{\omega}=A/K_{\omega}.
\]
The equivalence of \textup{(ii)} and \textup{(iii)} is exactly \Cref{prop:chain-vs-homological-visibility}(iii).
\end{proof}

\begin{theorem}[Character-blind pure-extension obstructions]\label{thm:character-blind-obstructions}
In the setting of \Cref{def:finite-nerve-presentation}, the following are equivalent:
\begin{enumerate}[label=(\roman*),nosep]
\item $K_{\omega}=0$;
\item $A_{\mathrm{vis}}^{\omega}=A$;
\item $\mathrm{ev}_{\omega}=0$;
\item $\omega$ lies in the subgroup
\[
\operatorname{Ext}^1(H_1(N(\mathcal U),\mathbb Z),A)\subseteq H^2(N(\mathcal U),A).
\]
\end{enumerate}
If these conditions hold and $\omega\neq 0$, then the chosen branch gerbe $\mathfrak G_v$ is non-neutral, while every character of $A$ is class-admissible for that branch. In particular, if $|\Vis_{p,s}(r)|=1$, then
\[
M,p\Vdash \Null^{\mathrm{glue}}_s(r).
\]
\end{theorem}

\begin{proof}
By \Cref{def:intrinsic-visible-quotient},
\[
K_{\omega}=\mathrm{Im}(\mathrm{ev}_{\omega}),
\]
so \textup{(i)} and \textup{(iii)} are equivalent. Since $A_{\mathrm{vis}}^{\omega}=A/K_{\omega}$, conditions \textup{(i)} and \textup{(ii)} are equivalent. By the universal coefficient sequence, \textup{(iii)} holds exactly when $\omega\in\ker(\eta_A)$, and this kernel is the stated $\operatorname{Ext}$-subgroup. Hence all four conditions are equivalent.

If they hold and $\omega\neq 0$, then \Cref{thm:cech-bridge-compatible-realizations}(iv) shows that the branch gerbe $\mathfrak G_v$ is non-neutral.
On the other hand, $\mathrm{ev}_{\omega}=0$ implies
\[
\chi\circ \mathrm{ev}_{\omega}=0
\]
for every $\chi\in\widehat A$, so \Cref{thm:intrinsic-visible-quotient} shows that every character is class-admissible.
\end{proof}

\begin{corollary}[Automatic blindness when $H_2$ vanishes]\label{cor:h2-vanishing-blindness}
Assume
\[
H_2(N(\mathcal U),\mathbb Z)=0.
\]
Then, for every obstruction class $\omega\in H^2(N(\mathcal U),A)$,
\[
\mathrm{ev}_{\omega}=0,
\qquad
K_{\omega}=0,
\qquad
A_{\mathrm{vis}}^{\omega}=A.
\]
Every nonzero branch obstruction class on such a nerve is therefore character-blind.
\end{corollary}

\begin{proof}
If $H_2(N(\mathcal U),\mathbb Z)=0$, then
\[
\mathrm{Hom}(H_2(N(\mathcal U),\mathbb Z),A)=0,
\]
so $\mathrm{ev}_{\omega}=0$ for every $\omega$. Apply \Cref{thm:character-blind-obstructions}.
\end{proof}

\begin{corollary}[No character-blind obstructions when $H_1$ is free]\label{cor:h1-free-no-blindness}
Assume that $H_1(N(\mathcal U),\mathbb Z)$ is free abelian. Then
\[
K_{\omega}=0
\quad\Longleftrightarrow\quad
\omega=0.
\]
Equivalently, every nonzero branch obstruction class has a nontrivial $H_2$-visible part. In particular, if $|\Vis_{p,s}(r)|=1$, every nontrivial gluing obstruction on such a nerve produces a proper intrinsic visible quotient.
\end{corollary}

\begin{proof}
If $H_1(N(\mathcal U),\mathbb Z)$ is free abelian, then
\[
\operatorname{Ext}^1(H_1(N(\mathcal U),\mathbb Z),A)=0.
\]
Hence the universal coefficient map $\eta_A$ is injective. By \Cref{thm:character-blind-obstructions},
\[
K_{\omega}=0
\quad\Longleftrightarrow\quad
\mathrm{ev}_{\omega}=0
\quad\Longleftrightarrow\quad
\eta_A(\omega)=0.
\]
Injectivity of $\eta_A$ therefore gives $K_{\omega}=0$ if and only if $\omega=0$. If $\omega\neq 0$, then $K_{\omega}\neq 0$, so $A_{\mathrm{vis}}^{\omega}=A/K_{\omega}$ is a proper quotient of the finite group $A$.
\end{proof}

\begin{example}[A character-blind obstruction on an $\mathbb{R}P^2$-type nerve]\label{ex:rp2-character-blind}
Let $\mathcal U$ be a finite cover whose Cech nerve is a finite triangulation of $\mathbb{R}P^2$, and take
\[
A=\mathbb Z/2\mathbb Z.
\]
Then
\[
H_2(\mathbb{R}P^2,\mathbb Z)=0,
\qquad
H_1(\mathbb{R}P^2,\mathbb Z)\cong \mathbb Z/2\mathbb Z
\]
\cite[Sec.~2.2]{Hatcher2002}. The universal coefficient sequence therefore yields
\[
H^2(\mathbb{R}P^2,\mathbb Z/2\mathbb Z)
\cong
\operatorname{Ext}^1(\mathbb Z/2\mathbb Z,\mathbb Z/2\mathbb Z)
\cong
\mathbb Z/2\mathbb Z
\]
\cite[Sec.~3.1]{Hatcher2002}. Let $\omega$ be the nonzero class. Then
\[
\mathrm{ev}_{\omega}=0,
\qquad
K_{\omega}=0,
\qquad
A_{\mathrm{vis}}^{\omega}=A,
\]
by \Cref{thm:character-blind-obstructions}. Thus no class-level quotient collapse occurs at all, yet the obstruction class is nonzero. Hence this branch is completely character-blind. If it is the unique visible branch of the admitted reference under consideration, then \Cref{thm:gerbe-null-semantics} yields a genuine gluing-level absence that is invisible to every character channel.
\end{example}

\begin{example}[A single-branch maximal-collapse obstruction on $S^2$]\label{ex:s2-single-branch-maximal-collapse}
Let $\mathcal U$ be a finite good cover of $S^2$ with nerve $N := N(\mathcal U)$, and take $A = \mathbb{Z}/2\mathbb{Z}$. Since
\[
H_1(N,\mathbb Z) = 0,
\qquad
H_2(N,\mathbb Z) \cong \mathbb Z
\]
\cite[Secs.~2.1, 2.2]{Hatcher2002}, the universal coefficient theorem gives
\[
H^2(N, \mathbb{Z}/2\mathbb{Z})
\cong
\mathrm{Hom}(\mathbb{Z}, \mathbb{Z}/2\mathbb{Z})
\cong
\mathbb{Z}/2\mathbb{Z}.
\]
Let $\omega$ be the nonzero class, and let $[S^2]$ denote the fundamental class in $H_2(N,\mathbb Z)$. Then
\[
\mathrm{ev}_\omega([S^2]) = 1 \in \mathbb{Z}/2\mathbb{Z},
\]
so the homological evaluation map is surjective. By \Cref{cor:maximal-collapse},
\[
K_\omega = A = \mathbb{Z}/2\mathbb{Z},
\qquad
A_{\mathrm{vis}}^\omega = 0.
\]
In the unique-branch case ($|\Vis_{p,s}(r)| = 1$), this gives an admitted reference in gluing-level absence for which the intrinsic visible quotient is trivial: every character channel is class-admissible, but the visible quotient collapses entirely. The full obstruction is $H_2$-visible and the strictification cost is $|K_\omega| = 2$. This is the opposite extreme from the $\mathbb{R}P^2$ blind case of \Cref{ex:rp2-character-blind}: there the obstruction is purely $\operatorname{Ext}$-type with maximal $A_{\mathrm{vis}}^\omega = A$, while here it is purely $H_2$-type with $A_{\mathrm{vis}}^\omega = 0$.
\end{example}

\begin{definition}[Intrinsic strictification code]\label{def:intrinsic-strictification-code}
An \emph{intrinsic strictification code} for the obstruction class $\omega$ consists of a finite set $R$ and an injective map
\[
\iota_{\omega}:A\hookrightarrow A_{\mathrm{vis}}^{\omega}\times R
\]
such that
\[
\mathrm{pr}_1\circ\iota_{\omega}=\pi_{\omega}.
\]
\end{definition}

\begin{corollary}[Intrinsic hidden-state lower bound]\label{cor:intrinsic-hidden-state-lower-bound}
Every intrinsic strictification code satisfies
\[
|R|\ge |K_{\omega}|.
\]
Equivalently,
\[
\log_2|R|\ge \log_2|K_{\omega}|.
\]
This lower bound is sharp.
\end{corollary}

\begin{proof}
By \Cref{def:intrinsic-visible-quotient}, the fibers of $\pi_{\omega}$ are exactly the cosets of $K_{\omega}$, so every fiber has cardinality $|K_{\omega}|$. The same fiber-counting argument as in \Cref{thm:minimal-information-cost-of-strictification} therefore shows that any injective code preserving the first projection must satisfy $|R|\ge |K_{\omega}|$.

For sharpness, take $R:=K_{\omega}$ and choose a set-theoretic section $s_{\omega}:A_{\mathrm{vis}}^{\omega}\to A$ of $\pi_{\omega}$. Then
\[
\iota_{\omega}(x):=\bigl(\pi_{\omega}(x),x-s_{\omega}(\pi_{\omega}(x))\bigr)
\]
is injective.
\end{proof}

\begin{theorem}[Prime decomposition of intrinsic visibility and budget]\label{thm:intrinsic-prime-decomposition}
Write the finite abelian group $A$ as a direct sum of its $p$-primary components:
\[
A=\bigoplus_p A_{(p)}.
\]
Let $\omega_p\in H^2(N(\mathcal U),A_{(p)})$ be the image of $\omega$ under the projection onto the $p$-primary summand. Then
\[
\mathrm{ev}_{\omega}=\bigoplus_p \mathrm{ev}_{\omega_p},
\qquad
K_{\omega}=\bigoplus_p K_{\omega,p},
\qquad
A_{\mathrm{vis}}^{\omega}\cong \bigoplus_p A_{(p)}/K_{\omega,p},
\]
where
\[
K_{\omega,p}:=\mathrm{Im}(\mathrm{ev}_{\omega_p}).
\]
Consequently,
\[
|K_{\omega}|=\prod_p |K_{\omega,p}|,
\qquad
\log_2|K_{\omega}|=\sum_p \log_2|K_{\omega,p}|.
\]
\end{theorem}

\begin{proof}
Both $H^2(N(\mathcal U),-)$ and $\mathrm{Hom}(H_2(N(\mathcal U),\mathbb Z),-)$ commute with finite direct sums, so the universal coefficient map is compatible with the primary decomposition of $A$. Hence
\[
\eta_A(\omega)=\bigoplus_p \eta_{A_{(p)}}(\omega_p),
\]
which is exactly the asserted decomposition of $\mathrm{ev}_{\omega}$.

Taking images gives
\[
K_{\omega}
=
\mathrm{Im}(\mathrm{ev}_{\omega})
=
\bigoplus_p \mathrm{Im}(\mathrm{ev}_{\omega_p})
=
\bigoplus_p K_{\omega,p},
\]
by \Cref{def:intrinsic-visible-quotient}. The quotient decomposition follows immediately:
\[
A_{\mathrm{vis}}^{\omega}
=
A/K_{\omega}
\cong
\bigoplus_p A_{(p)}/K_{\omega,p}.
\]
Since finite direct sums correspond to products of cardinalities, the budget formula is immediate.
\end{proof}

\begin{corollary}[Intrinsic budget versus strict budget]\label{cor:strict-intrinsic-budget-comparison}
For every finite obstruction presentation $\alpha$ of the class $\omega$,
\[
|K_{\omega}|\le |H_\alpha|,
\qquad
\log_2|K_{\omega}|\le \log_2|H_\alpha|.
\]
Hence the intrinsic hidden-state budget never exceeds the strict presentation-level budget. Equality holds if and only if the canonical surjection
\[
A_{\mathrm{vis}}^{\omega}\twoheadrightarrow A/H_\alpha
\]
is an isomorphism.
\end{corollary}

\begin{proof}
By \Cref{cor:strict-to-intrinsic-visible}, there is a surjective homomorphism
\[
A_{\mathrm{vis}}^{\omega}=A/K_{\omega}\twoheadrightarrow A/H_\alpha.
\]
For finite groups this implies
\[
|A|/|K_{\omega}|\ge |A|/|H_\alpha|,
\]
hence $|K_{\omega}|\le |H_\alpha|$. The logarithmic inequality is immediate. Equality of cardinalities holds exactly when the displayed surjection is bijective.
\end{proof}

\begin{corollary}[Maximal collapse criterion]\label{cor:maximal-collapse}
In the setting of \Cref{def:finite-nerve-presentation}, the following are equivalent:
\begin{enumerate}[label=(\roman*),nosep]
\item $\mathrm{ev}_{\omega}:H_2(N(\mathcal U),\mathbb Z)\to A$ is surjective;
\item $K_{\omega}=A$;
\item $A_{\mathrm{vis}}^{\omega}=0$;
\item the intrinsic strictification budget is maximal:
\[
\log_2|K_{\omega}|=\log_2|A|.
\]
\end{enumerate}
\end{corollary}

\begin{proof}
By \Cref{def:intrinsic-visible-quotient}, $K_{\omega}=\mathrm{Im}(\mathrm{ev}_{\omega})$, so \textup{(i)} and \textup{(ii)} are equivalent. Condition \textup{(ii)} is equivalent to \textup{(iii)} by the definition of $A_{\mathrm{vis}}^{\omega}=A/K_{\omega}$. For finite groups, \textup{(ii)} is also equivalent to $|K_{\omega}|=|A|$, which is exactly \textup{(iv)} after taking $\log_2$.
\end{proof}

\subsection{Aggregating visibility across branches}

The preceding results analyze one visible branch at a time. When an admitted reference carries several visible branches, there are two distinct ways of aggregating the branchwise visible data: one may allow the observer to know which branch has been reached, or one may require the same observation channel to remain admissible on every branch simultaneously. These two requirements lead to different quotients.

\begin{definition}[Branchwise intrinsic data]\label{def:branchwise-intrinsic-data}
Fix an admitted reference $r:\mathsf{Ref}_s$ at a state $p$, and assume:
\begin{enumerate}[label=(\roman*),nosep]
\item the hypotheses of \Cref{def:cech-gluing-obstruction} hold;
\item the set $\Vis_{p,s}(r)$ is finite and nonempty;
\item every component gerbe $\mathfrak G_v$ is banded by the same finite abelian group $A$ and admits a finite nerve presentation of its branch obstruction class $\omega_v$.
\end{enumerate}
For each visible branch $v\in \Vis_{p,s}(r)$, write
\[
K_v:=K_{\omega_v},
\qquad
A_{\mathrm{vis}}^v:=A_{\mathrm{vis}}^{\omega_v}=A/K_v,
\qquad
\pi_v:A\twoheadrightarrow A_{\mathrm{vis}}^v.
\]
\end{definition}

\begin{theorem}[Branch-sensitive visible quotient]\label{thm:branch-sensitive-visible-quotient}
In the setting of \Cref{def:branchwise-intrinsic-data}, consider the diagonal map
\[
\Delta_{\mathrm{br}}:A\longrightarrow \prod_{v\in \Vis_{p,s}(r)} A_{\mathrm{vis}}^v,
\qquad
x\longmapsto (\pi_v(x))_v.
\]
Then
\[
\ker(\Delta_{\mathrm{br}})=\bigcap_{v\in \Vis_{p,s}(r)} K_v.
\]
Consequently, the quotient
\[
A_{\mathrm{vis}}^{\mathrm{br}}(r):=A\Big/\bigcap_{v\in \Vis_{p,s}(r)} K_v
\]
is canonically isomorphic to $\mathrm{Im}(\Delta_{\mathrm{br}})$. It is the maximal quotient of $A$ that remains visible once the branch label itself is allowed as part of the semantic context.
\end{theorem}

\begin{proof}
For $x\in A$ one has $\Delta_{\mathrm{br}}(x)=0$ if and only if $\pi_v(x)=0$ for every visible branch $v$, which is equivalent to $x\in K_v$ for every $v$. Hence
\[
\ker(\Delta_{\mathrm{br}})=\bigcap_{v\in \Vis_{p,s}(r)} K_v.
\]
The first isomorphism theorem therefore yields
\[
\mathrm{Im}(\Delta_{\mathrm{br}})
\cong
A\Big/\bigcap_{v\in \Vis_{p,s}(r)} K_v.
\]
The final interpretation is immediate: once the visible branch is known, two elements of $A$ remain semantically indistinguishable exactly when they are indistinguishable on every individual branch quotient.
\end{proof}

\begin{theorem}[Branch-uniform visible quotient]\label{thm:branch-uniform-visible-quotient}
In the setting of \Cref{def:branchwise-intrinsic-data}, let
\[
X_v:=\{\chi\in \widehat A\mid \chi_*(\omega_v)=0\}
\]
for each visible branch $v$, and define the uniformly admissible characters by
\[
X^{\mathrm{unif}}(r):=\bigcap_{v\in \Vis_{p,s}(r)} X_v.
\]
Then
\[
X^{\mathrm{unif}}(r)=\mathrm{Ann}\Bigl(\sum_{v\in \Vis_{p,s}(r)} K_v\Bigr).
\]
Consequently, the quotient
\[
A_{\mathrm{vis}}^{\mathrm{unif}}(r):=A\Big/\sum_{v\in \Vis_{p,s}(r)} K_v
\]
is exactly the quotient detected by the characters that are admissible on every visible branch simultaneously.
\end{theorem}

\begin{proof}
By \Cref{thm:intrinsic-visible-quotient}, one has
\[
X_v=\mathrm{Ann}(K_v)
\]
for every $v$. Therefore
\[
X^{\mathrm{unif}}(r)
=
\bigcap_{v\in \Vis_{p,s}(r)} \mathrm{Ann}(K_v)
=
\mathrm{Ann}\Bigl(\sum_{v\in \Vis_{p,s}(r)} K_v\Bigr),
\]
where the last identity is the standard annihilator formula for finite abelian groups \cite[Chap.~5]{DummitFoote2004}. Finite abelian duality then gives
\[
\bigcap_{\chi\in X^{\mathrm{unif}}(r)} \ker\chi
=
\sum_{v\in \Vis_{p,s}(r)} K_v.
\]
Hence the quotient detected by all uniformly admissible characters is precisely
\[
A\Big/\sum_{v\in \Vis_{p,s}(r)} K_v.
\qedhere
\]
\end{proof}

\begin{definition}[Branch comparison map and resolution kernel]\label{def:branch-comparison-map}
In the setting of \Cref{def:branchwise-intrinsic-data}, let
\[
\pi_{\mathrm{br}}:A\twoheadrightarrow A_{\mathrm{vis}}^{\mathrm{br}}(r)
\qquad\text{and}\qquad
\pi_{\mathrm{unif}}:A\twoheadrightarrow A_{\mathrm{vis}}^{\mathrm{unif}}(r)
\]
be the canonical quotient maps. Since
\[
\bigcap_{v\in \Vis_{p,s}(r)} K_v
\subseteq
\sum_{v\in \Vis_{p,s}(r)} K_v,
\]
there is a unique surjective homomorphism
\[
\beta_r:A_{\mathrm{vis}}^{\mathrm{br}}(r)\twoheadrightarrow A_{\mathrm{vis}}^{\mathrm{unif}}(r)
\]
such that
\[
\beta_r\circ \pi_{\mathrm{br}}=\pi_{\mathrm{unif}}.
\]
Its kernel is denoted
\[
B_{\mathrm{res}}(r):=
\ker(\beta_r)
\cong
\Bigl(\sum_{v\in \Vis_{p,s}(r)} K_v\Bigr)\Big/\Bigl(\bigcap_{v\in \Vis_{p,s}(r)} K_v\Bigr)
\]
and called the \emph{branch-resolution kernel}.
\end{definition}

\begin{theorem}[Exact comparison of branch-sensitive and branch-uniform visibility]\label{thm:branch-comparison-exact-sequence}
In the setting of \Cref{def:branch-comparison-map}, there is a short exact sequence
\[
0\longrightarrow B_{\mathrm{res}}(r)\longrightarrow A_{\mathrm{vis}}^{\mathrm{br}}(r)\xrightarrow{\ \beta_r\ } A_{\mathrm{vis}}^{\mathrm{unif}}(r)\longrightarrow 0.
\]
Equivalently,
\[
\ker(\beta_r)
\cong
\Bigl(\sum_{v\in \Vis_{p,s}(r)} K_v\Bigr)\Big/\Bigl(\bigcap_{v\in \Vis_{p,s}(r)} K_v\Bigr).
\]
Thus the precise gap between branch-sensitive and branch-uniform visibility is measured by the lattice interval from $\bigcap_v K_v$ to $\sum_v K_v$.
\end{theorem}

\begin{proof}
Existence and uniqueness of $\beta_r$ follow from the universal property of the quotient because
\[
\bigcap_v K_v\subseteq \sum_v K_v.
\]
Surjectivity is immediate from surjectivity of $\pi_{\mathrm{unif}}$. By the third isomorphism theorem,
\[
\ker(\beta_r)
\cong
\Bigl(\sum_v K_v\Bigr)\Big/\Bigl(\bigcap_v K_v\Bigr),
\]
which gives the claimed exact sequence.
\end{proof}

\begin{corollary}[Criterion for the absence of branch gain]\label{cor:no-branch-gain-criterion}
In the setting of \Cref{def:branch-comparison-map}, the following are equivalent:
\begin{enumerate}[label=(\roman*),nosep]
\item $\beta_r$ is an isomorphism;
\item $B_{\mathrm{res}}(r)=0$;
\item
\[
\bigcap_{v\in \Vis_{p,s}(r)} K_v
=
\sum_{v\in \Vis_{p,s}(r)} K_v;
\]
\item the hidden subgroups are branch-constant:
\[
K_v=K_w
\qquad
(\forall v,w\in \Vis_{p,s}(r)).
\]
\end{enumerate}
Hence branch sensitivity yields no additional visible information exactly when all visible branches have the same intrinsic hidden subgroup.
\end{corollary}

\begin{proof}
The equivalence of \textup{(i)} and \textup{(ii)} is exactness of the sequence in \Cref{thm:branch-comparison-exact-sequence}. Condition \textup{(ii)} is equivalent to \textup{(iii)} by the quotient description of $B_{\mathrm{res}}(r)$.

If \textup{(iii)} holds, write the common subgroup as $H$. Then
\[
H=\bigcap_v K_v\subseteq K_v\subseteq \sum_v K_v=H
\]
for every visible branch $v$, so each $K_v=H$, proving \textup{(iv)}. Conversely, if all $K_v$ are equal to one subgroup $H$, then clearly
\[
\bigcap_v K_v=H=\sum_v K_v.
\qedhere
\]
\end{proof}

\begin{corollary}[Dual character description of the branch-resolution kernel]\label{cor:branch-resolution-dual}
In the setting of \Cref{def:branch-comparison-map},
\[
\widehat{A_{\mathrm{vis}}^{\mathrm{br}}(r)}
\cong
\sum_{v\in \Vis_{p,s}(r)} X_v,
\qquad
\widehat{A_{\mathrm{vis}}^{\mathrm{unif}}(r)}
\cong
\bigcap_{v\in \Vis_{p,s}(r)} X_v.
\]
Moreover,
\[
\widehat{B_{\mathrm{res}}(r)}
\cong
\Bigl(\sum_{v\in \Vis_{p,s}(r)} X_v\Bigr)\Big/\Bigl(\bigcap_{v\in \Vis_{p,s}(r)} X_v\Bigr).
\]
Thus the character channels created by branch resolution are exactly the channels that belong to the subgroup generated by the branchwise admissible characters but fail to be uniformly admissible.
\end{corollary}

\begin{proof}
By \Cref{thm:intrinsic-visible-quotient}, one has
\[
X_v=\mathrm{Ann}(K_v)
\]
for every visible branch $v$. Finite abelian duality then gives
\[
\widehat{A_{\mathrm{vis}}^{\mathrm{br}}(r)}
\cong
\mathrm{Ann}\Bigl(\bigcap_v K_v\Bigr)
=
\sum_v \mathrm{Ann}(K_v)
=
\sum_v X_v
\]
and, by \Cref{thm:branch-uniform-visible-quotient},
\[
\widehat{A_{\mathrm{vis}}^{\mathrm{unif}}(r)}
\cong
\bigcap_v X_v.
\]
Dualizing the short exact sequence of \Cref{thm:branch-comparison-exact-sequence} yields an exact sequence
\[
0\longrightarrow \widehat{A_{\mathrm{vis}}^{\mathrm{unif}}(r)}
\longrightarrow \widehat{A_{\mathrm{vis}}^{\mathrm{br}}(r)}
\longrightarrow \widehat{B_{\mathrm{res}}(r)}
\longrightarrow 0,
\]
and substituting the two identifications above gives the claimed quotient description of $\widehat{B_{\mathrm{res}}(r)}$.
\end{proof}

\begin{definition}[Branch-resolution code]\label{def:branch-resolution-code}
In the setting of \Cref{def:branch-comparison-map}, a \emph{branch-resolution code} consists of a finite set $R$ and an injective map
\[
\iota_{\mathrm{res}}:A_{\mathrm{vis}}^{\mathrm{br}}(r)\hookrightarrow A_{\mathrm{vis}}^{\mathrm{unif}}(r)\times R
\]
such that
\[
\mathrm{pr}_1\circ \iota_{\mathrm{res}}=\beta_r.
\]
\end{definition}

\begin{theorem}[Exact branch-resolution budget]\label{thm:branch-resolution-budget}
Every branch-resolution code satisfies
\[
|R|\ge |B_{\mathrm{res}}(r)|
=
\left|\Bigl(\sum_v K_v\Bigr)\Big/\Bigl(\bigcap_v K_v\Bigr)\right|.
\]
Equivalently,
\[
\log_2|R|
\ge
\log_2|B_{\mathrm{res}}(r)|.
\]
This bound is sharp.
\end{theorem}

\begin{proof}
Since $\mathrm{pr}_1\circ \iota_{\mathrm{res}}=\beta_r$, every fiber of $\beta_r$ injects into $R$. By \Cref{thm:branch-comparison-exact-sequence}, each fiber of $\beta_r$ is a coset of $B_{\mathrm{res}}(r)$ and therefore has cardinality $|B_{\mathrm{res}}(r)|$. Hence
\[
|R|\ge |B_{\mathrm{res}}(r)|.
\]

For sharpness, choose a finite set $R$ with
\[
|R|=|B_{\mathrm{res}}(r)|.
\]
Choose a set-theoretic section
\[
s:A_{\mathrm{vis}}^{\mathrm{unif}}(r)\to A_{\mathrm{vis}}^{\mathrm{br}}(r)
\]
of $\beta_r$, and choose an injection
\[
j:B_{\mathrm{res}}(r)\hookrightarrow R.
\]
Define
\[
\iota_{\mathrm{res}}(x):=
\bigl(\beta_r(x),j(x-s(\beta_r(x)))\bigr).
\]
Since $x-s(\beta_r(x))\in \ker(\beta_r)=B_{\mathrm{res}}(r)$, the map is well defined. If
\[
\iota_{\mathrm{res}}(x)=\iota_{\mathrm{res}}(y),
\]
then $\beta_r(x)=\beta_r(y)$ and
\[
j(x-s(\beta_r(x)))=j(y-s(\beta_r(y))).
\]
Injectivity of $j$ gives $x=y$, so $\iota_{\mathrm{res}}$ is injective and the bound is attained.
\end{proof}

\begin{corollary}[Factorization of uniform hidden information]\label{cor:uniform-hidden-budget-factorization}
In the setting of \Cref{def:branch-comparison-map},
\[
\left|\sum_{v\in \Vis_{p,s}(r)} K_v\right|
=
|B_{\mathrm{res}}(r)|
\cdot
\left|\bigcap_{v\in \Vis_{p,s}(r)} K_v\right|.
\]
Equivalently,
\[
\log_2\left|\sum_v K_v\right|
=
\log_2|B_{\mathrm{res}}(r)|
+
\log_2\left|\bigcap_v K_v\right|.
\]
Thus the hidden-state budget behind uniform visibility factors exactly into a branch-resolution term and a residual within-branch hidden term.
\end{corollary}

\begin{proof}
By \Cref{def:branch-comparison-map},
\[
|B_{\mathrm{res}}(r)|
=
\left|\Bigl(\sum_v K_v\Bigr)\Big/\Bigl(\bigcap_v K_v\Bigr)\right|
=
\frac{|\sum_v K_v|}{|\bigcap_v K_v|}.
\]
Rearranging gives the first displayed identity, and the second follows after taking base-$2$ logarithms.
\end{proof}

\begin{definition}[Branch-preserving strictification code]\label{def:branch-preserving-strictification-code}
In the setting of \Cref{def:branchwise-intrinsic-data}, a \emph{branch-preserving strictification code} consists of a finite set $R$ and, for each visible branch $v\in \Vis_{p,s}(r)$, an injective map
\[
\iota_v:A\hookrightarrow A_{\mathrm{vis}}^v\times R
\]
such that
\[
\mathrm{pr}_1\circ \iota_v=\pi_v.
\]
\end{definition}

\begin{theorem}[Multi-branch strictification budget]\label{thm:multi-branch-strictification-budget}
Every branch-preserving strictification code satisfies
\[
|R|\ge \max_{v\in \Vis_{p,s}(r)} |K_v|.
\]
Equivalently,
\[
\log_2|R|\ge \max_{v\in \Vis_{p,s}(r)} \log_2|K_v|.
\]
This bound is sharp.
\end{theorem}

\begin{proof}
Fix a visible branch $v$. Since $\mathrm{pr}_1\circ \iota_v=\pi_v$, every fiber of $\pi_v$ injects into $R$. By \Cref{def:intrinsic-visible-quotient}, every fiber of $\pi_v$ is a coset of $K_v$ and hence has cardinality $|K_v|$. Therefore $|R|\ge |K_v|$ for every $v$, which proves the lower bound.

For sharpness, choose a finite set $R$ with
\[
|R|=\max_{v\in \Vis_{p,s}(r)} |K_v|.
\]
For each $v$, choose an injection
\[
j_v:K_v\hookrightarrow R
\]
and a set-theoretic section
\[
s_v:A_{\mathrm{vis}}^v\to A
\]
of $\pi_v$. Define
\[
\iota_v(x):=\bigl(\pi_v(x),j_v(x-s_v(\pi_v(x)))\bigr).
\]
Since $x-s_v(\pi_v(x))\in K_v$, the formula is well defined. If $\iota_v(x)=\iota_v(y)$, then $\pi_v(x)=\pi_v(y)$ and
\[
j_v(x-s_v(\pi_v(x)))=j_v(y-s_v(\pi_v(y))).
\]
Injectivity of $j_v$ gives $x=y$. Thus each $\iota_v$ is injective, so the bound is attained.
\end{proof}

\begin{corollary}[Single-branch reduction]\label{cor:single-branch-reduction}
If
\[
\Vis_{p,s}(r)=\{v\},
\]
then
\[
A_{\mathrm{vis}}^{\mathrm{br}}(r)\cong A_{\mathrm{vis}}^{\mathrm{unif}}(r)\cong A_{\mathrm{vis}}^v
\]
and the optimal branch-preserving strictification budget is exactly $|K_v|$.
\end{corollary}

\begin{proof}
If there is only one visible branch, then
\[
\bigcap_{u\in \Vis_{p,s}(r)} K_u
=
\sum_{u\in \Vis_{p,s}(r)} K_u
=
K_v.
\]
The quotient statements follow from \Cref{thm:branch-sensitive-visible-quotient,thm:branch-uniform-visible-quotient}, and the budget statement is \Cref{thm:multi-branch-strictification-budget}.
\end{proof}

\begin{example}[A two-branch sphere example separating the aggregate invariants]\label{ex:two-branch-sphere-separation}
Let $\mathcal U$ be a finite good cover of $S^2$ with nerve
\[
N:=N(\mathcal U),
\]
and take
\[
A=(\mathbb Z/2\mathbb Z)^2.
\]
Since
\[
H_1(N,\mathbb Z)=0,
\qquad
H_2(N,\mathbb Z)\cong \mathbb Z
\]
\cite[Secs.~2.1, 2.2]{Hatcher2002}, the universal coefficient theorem gives
\[
H^2(N,A)\cong \mathrm{Hom}(H_2(N,\mathbb Z),A)\cong A.
\]
Choose classes
\[
\omega_1,\omega_2\in H^2(N,A)
\]
corresponding to the two basis vectors $(1,0)$ and $(0,1)$, and let
\[
\mathfrak G_{v_1},
\qquad
\mathfrak G_{v_2}
\]
be banded $A$-gerbes representing those classes \cite[Chap.~IV]{Giraud1971}, \cite[Tag 06NY]{StacksProject}. Put
\[
\mathfrak L_r:=\mathfrak G_{v_1}\sqcup \mathfrak G_{v_2}
\]
and define the local object functor on $\mathcal C_a$ by
\[
F_{p,s}|_{\mathcal C_a}:=\pi_0(\mathfrak L_r).
\]
Since $\mathfrak L_r$ is already a stack, we may take
\[
\mathfrak P_r:=\mathfrak L_r
\]
as the realization prestack of an admitted reference $r$, with
\[
\pi_0^{\mathrm{pre}}(\mathfrak P_r)=\pi_0(\mathfrak L_r)=F_{p,s}|_{\mathcal C_a}.
\]
Then
\[
\Vis_{p,s}(r)=\{v_1,v_2\},
\]
and
\[
M,p\Vdash \CompSec_s(r)\land \neg \Sec_s(r),
\]
because each gerbe is locally nonempty while both are non-neutral.

If $[S^2]$ denotes the fundamental class in $H_2(N,\mathbb Z)$, then
\[
\mathrm{ev}_{\omega_1}([S^2])=(1,0),
\qquad
\mathrm{ev}_{\omega_2}([S^2])=(0,1).
\]
Hence
\[
K_{v_1}=\langle(1,0)\rangle,
\qquad
K_{v_2}=\langle(0,1)\rangle,
\qquad
A_{\mathrm{vis}}^{v_1}\cong A_{\mathrm{vis}}^{v_2}\cong \mathbb Z/2\mathbb Z.
\]
Therefore
\[
\bigcap_{v} K_v=0,
\qquad
\sum_{v} K_v=A,
\qquad
\max_v |K_v|=2.
\]
By \Cref{thm:branch-sensitive-visible-quotient,thm:branch-uniform-visible-quotient,thm:multi-branch-strictification-budget},
\[
A_{\mathrm{vis}}^{\mathrm{br}}(r)\cong A,
\qquad
A_{\mathrm{vis}}^{\mathrm{unif}}(r)=0,
\qquad
\log_2\!\bigl(\max_v |K_v|\bigr)=1.
\]
Moreover,
\[
B_{\mathrm{res}}(r)\cong A,
\qquad
\log_2|B_{\mathrm{res}}(r)|=2,
\qquad
\left|\bigcap_v K_v\right|=1.
\]
Thus each branch individually hides one coordinate of $A$, but the hidden coordinates are transverse: once the branch label is known, nothing remains hidden across all branches, whereas any single character channel admissible on both branches must be trivial. In this example the entire uniform hidden budget is exhausted by branch resolution itself.
\end{example}

\subsection{Connection to sheaf-theoretic contextuality}

The unique-branch case $|\Vis_{p,s}(r)|=1$ deserves separate attention because it connects the present theory to the sheaf-theoretic analysis of contextuality. The comparison can be stated formally once the local object functor is specialized to the support presheaf of an Abramsky--Brandenburger measurement scenario.

\begin{theorem}[Unique-branch comparison with the no-global-section picture]\label{thm:unique-branch-contextuality-comparison}
Let $r:\mathsf{Ref}_s$ be admitted at $p$, write $a=[r]_p$, and assume:
\begin{enumerate}[label=(\roman*),nosep]
\item the slice site $\mathcal C_a$ is the measurement-cover site of a finite Abramsky--Brandenburger scenario, and the local object functor $F_{p,s}|_{\mathcal C_a}$ is identified with the support presheaf $S_e$ of an empirical model $e$ \cite{AbramskyBrandenburger2011};
\item $r$ admits a realization prestack that is globally conservative at $a$;
\item $|\Vis_{p,s}(r)|=1$, with unique visible branch $v$;
\item the branch obstruction class
\[
\omega:=\omega_v
\]
admits a finite nerve presentation with band $A$.
\end{enumerate}
Then:
\begin{enumerate}[label=(\roman*),nosep]
\item
\[
M,p\Vdash \CompSec_s(r)
\quad\Longleftrightarrow\quad
S_e\ \text{has a compatible family};
\]
\item
\[
M,p\Vdash \Sec_s(r)
\quad\Longleftrightarrow\quad
S_e\ \text{has a global section};
\]
\item consequently,
\[
M,p\Vdash \Null^{\mathrm{glue}}_s(r)
\quad\Longleftrightarrow\quad
S_e\ \text{has a compatible family but no global section};
\]
\item
\[
A_{\mathrm{vis}}^{\omega}
\cong
A_{\mathrm{vis}}^{\mathrm{br}}(r)
\cong
A_{\mathrm{vis}}^{\mathrm{unif}}(r),
\]
and this quotient is exactly the quotient of $A$ detected by the class-admissible characters;
\item if $\omega\neq 0$, then
\[
A_{\mathrm{vis}}^{\omega}=A
\quad\Longleftrightarrow\quad
K_{\omega}=0
\quad\Longleftrightarrow\quad
\omega\in \operatorname{Ext}^1(H_1(N(\mathcal U),\mathbb Z),A).
\]
Hence the no-global-section obstruction can remain nontrivial while being invisible to every admissible character channel precisely when the surviving class is pure $\operatorname{Ext}$-type.
\end{enumerate}
\end{theorem}

\begin{proof}
Under the identification $F_{p,s}|_{\mathcal C_a}\cong S_e$, a compatible family in the support presheaf is exactly a compatible local family in the sense of \Cref{def:existence-levels}(ii), and a global section of $S_e$ is exactly an element of $F_{p,s}(a)$. This proves \textup{(i)} and \textup{(ii)}.

Statement \textup{(iii)} is then just the definition of $\Null^{\mathrm{glue}}_s(r)$. Since there is only one visible branch, \Cref{cor:single-branch-reduction} gives the quotient identifications in \textup{(iv)}, and \Cref{cor:intrinsic-character-detection} identifies that quotient with the one detected by the class-admissible characters. Finally, \textup{(v)} is exactly \Cref{thm:character-blind-obstructions} rewritten for the unique branch $\omega$.
\end{proof}

\begin{remark}[Car\`u-type blind cases as pure $\operatorname{Ext}$-residuals]\label{rem:caru-ext-residual}
\Cref{thm:unique-branch-contextuality-comparison}(v) turns the comparison with \cite{Caru2017} into a theorem: the failure of cohomological completeness in the unique-branch case is not an external phenomenon but exactly the case in which the surviving obstruction class lies entirely in the $\operatorname{Ext}$-summand of the universal coefficient sequence.
\end{remark}
